\subsection{由集合生成的拟伽罗瓦扩张}

命题 4.——设 \((\mathbf{N}_i)_{i\in I}\) 是包含在 \(\Omega\) 中的 K 的一族拟伽罗瓦扩张。令 \(\mathbf{N} = \bigcap_{i\in I}\mathbf{N}_i\) 且 \(\mathbf{M} = \mathbf{K}\left(\bigcup_{i\in I}\mathbf{N}_i\right)\) ；则扩张 N 与 M

在 K 上都是拟伽罗瓦的。

设 \(u\) 是 \(\Omega\) 的一个 K-自同构。我们有 \(u(\mathbf{N}_i) = \mathbf{N}_i\) （对所有 \(i \in I\) ），由此得到等式 \(u(\mathbf{N}) = \mathbf{N}\) 与 \(u(\mathbf{M}) = \mathbf{M}\) ；于是命题由推论 1（V，第 54 页）得出。

设 A 是由 \(\Omega\) 的元素组成的一个集合，B 是由 \(\Omega\) 中与 A 的元素共轭的元素组成的集合；换句话说，若 G 是 \(\Omega\) 的 K-自同构的群，则我们有 \(\mathbf{B} = \bigcup_{u \in G} u(\mathbf{A})\) 。对每个 \(u \in G\) 我们有

\(u(B) = B\) ，从而 \(u(K(B)) = K(B)\) 。因此 \(K(B)\) 是 K 的一个包含 A 的拟伽罗瓦扩张，且显然 K 的每个包含 A 的拟伽罗瓦扩张 N 都包含 B，从而包含 \(K(B)\) 。我们称 \(K(B)\) 为由 A 生成的拟伽罗瓦扩张。这一定义特别适用于 A 是 K 的扩张的情形。

下述命题直接由前面的注得到。

命题 5.——设 E 是包含在 Ω 中的 K 的一个扩张，N 是由 E 生成的拟伽罗瓦扩张。若 A 是 Ω 的一个子集使得 E = K(A)，则 N = K(B)，其中 B 是 Ω 中与 A 的某个元素共轭的元素组成的集合。

推论 1.——若 E 是 K 的有限次扩张，则由 E 生成的 K 的拟伽罗瓦扩张 N 在 K 上是有限次的。

因为我们有 \(E = K(A)\) ，其中 A 是有限的，从而 A 的元素的共轭元所成的集 B 是有限的，于是推论由定理 2（V，第 18 页）得出。

推论 2.——K 的每个拟伽罗瓦扩张 N 都是 N 的在 K 上有限的拟伽罗瓦子扩张之并。

设 \(x \in \mathbb{N}\) ，并设 \(\mathbf{N}_x\) 是 \(\mathbf{K}\) 的一个拟伽罗瓦扩张（由 \(\{x\}\) 生成的）。由于 \(\mathbf{K}(x)\) 在 \(\mathbf{K}\) 上是有限次的，\(\mathbf{N}_x\) 亦然（推论 1），且我们有 \(x \in \mathbf{N}_x\) .

